\documentclass[10pt,a4paper]{amsart}
\usepackage{amsmath,amssymb,amsthm,graphicx}
% Retained minimal preamble for the Stacks Project chapter more-morphisms.tex (More on Morphisms).
% Verbatim excerpt of the chapter preamble: Xy-pic setup (native preamble lines 14-18) and the native
% theorem declarations with the mathematical macros (native preamble lines 154-217). Copied mechanically;
% nothing added, removed or redefined. The wrapper class and the amsmath/amssymb/amsthm/graphicx base
% packages are supplied by the assembler, so the class/externaldocument part of the native preamble is omitted.
% For commutative diagrams we use Xy-pic
\usepackage[all]{xy}

% We use 2cell for 2-commutative diagrams.
\xyoption{2cell}
% Theorem environments.
%
\theoremstyle{plain}
\newtheorem{theorem}[subsection]{Theorem}
\newtheorem{proposition}[subsection]{Proposition}
\newtheorem{lemma}[subsection]{Lemma}

\theoremstyle{definition}
\newtheorem{definition}[subsection]{Definition}
\newtheorem{example}[subsection]{Example}
\newtheorem{exercise}[subsection]{Exercise}
\newtheorem{situation}[subsection]{Situation}

\theoremstyle{remark}
\newtheorem{remark}[subsection]{Remark}
\newtheorem{remarks}[subsection]{Remarks}

\numberwithin{equation}{subsection}

% Macros
%
\def\lim{\mathop{\mathrm{lim}}\nolimits}
\def\colim{\mathop{\mathrm{colim}}\nolimits}
\def\Spec{\mathop{\mathrm{Spec}}}
\def\Hom{\mathop{\mathrm{Hom}}\nolimits}
\def\Ext{\mathop{\mathrm{Ext}}\nolimits}
\def\SheafHom{\mathop{\mathcal{H}\!\mathit{om}}\nolimits}
\def\SheafExt{\mathop{\mathcal{E}\!\mathit{xt}}\nolimits}
\def\Sch{\mathit{Sch}}
\def\Mor{\mathop{\mathrm{Mor}}\nolimits}
\def\Ob{\mathop{\mathrm{Ob}}\nolimits}
\def\Sh{\mathop{\mathit{Sh}}\nolimits}
\def\NL{\mathop{N\!L}\nolimits}
\def\CH{\mathop{\mathrm{CH}}\nolimits}
\def\proetale{{pro\text{-}\acute{e}tale}}
\def\etale{{\acute{e}tale}}
\def\QCoh{\mathit{QCoh}}
\def\Ker{\mathop{\mathrm{Ker}}}
\def\Im{\mathop{\mathrm{Im}}}
\def\Coker{\mathop{\mathrm{Coker}}}
\def\Coim{\mathop{\mathrm{Coim}}}

% Boxtimes
%
\DeclareMathSymbol{\boxtimes}{\mathbin}{AMSa}{"02}

%
% Macros for moduli stacks/spaces
%
\def\QCohstack{\mathcal{QC}\!\mathit{oh}}
\def\Cohstack{\mathcal{C}\!\mathit{oh}}
\def\Spacesstack{\mathcal{S}\!\mathit{paces}}
\def\Quotfunctor{\mathrm{Quot}}
\def\Hilbfunctor{\mathrm{Hilb}}
\def\Curvesstack{\mathcal{C}\!\mathit{urves}}
\def\Polarizedstack{\mathcal{P}\!\mathit{olarized}}
\def\Complexesstack{\mathcal{C}\!\mathit{omplexes}}
% \Pic is the operator that assigns to X its picard group, usage \Pic(X)
% \Picardstack_{X/B} denotes the Picard stack of X over B
% \Picardfunctor_{X/B} denotes the Picard functor of X over B
\def\Pic{\mathop{\mathrm{Pic}}\nolimits}
\def\Picardstack{\mathcal{P}\!\mathit{ic}}
\def\Picardfunctor{\mathrm{Pic}}
\def\Deformationcategory{\mathcal{D}\!\mathit{ef}}
\newenvironment{reference}{}{}
\newenvironment{slogan}{}{}
\newenvironment{history}{}{}
\usepackage{cleveref}
\crefname{originallemma}{original lemma}{original lemmas}
\Crefname{originallemma}{Original lemma}{Original lemmas}
\makeatletter
\newlabel{lemma-NL-formally-smooth}{{lemma-NL-formally-smooth}{source}{}{sourcepointer-1}{}}
\newlabel{lemma-NL-formally-smooth@cref}{{[originallemma][0][]lemma-NL-formally-smooth}{[1][0][]source}{}{}{}}
\makeatother
\crefname{originallemma}{original lemma}{original lemmas}
\Crefname{originallemma}{Original lemma}{Original lemmas}
\makeatletter
\newlabel{lemma-NL-immersion}{{lemma-NL-immersion}{source}{}{sourcepointer-2}{}}
\newlabel{lemma-NL-immersion@cref}{{[originallemma][0][]lemma-NL-immersion}{[1][0][]source}{}{}{}}
\makeatother
\crefname{originallemma}{original lemma}{original lemmas}
\Crefname{originallemma}{Original lemma}{Original lemmas}
\makeatletter
\newlabel{lemma-smooth-formally-smooth}{{lemma-smooth-formally-smooth}{source}{}{sourcepointer-3}{}}
\newlabel{lemma-smooth-formally-smooth@cref}{{[originallemma][0][]lemma-smooth-formally-smooth}{[1][0][]source}{}{}{}}
\makeatother
\crefname{originallemma}{original lemma}{original lemmas}
\Crefname{originallemma}{Original lemma}{Original lemmas}
\makeatletter
\newlabel{modules-lemma-exact-sequence-NL-ringed-topoi}{{0E1Z}{source}{}{sourcepointer-4}{}}
\newlabel{modules-lemma-exact-sequence-NL-ringed-topoi@cref}{{[originallemma][0][]0E1Z}{[1][0][]source}{}{}{}}
\makeatother
\crefname{originallemma}{original lemma}{original lemmas}
\Crefname{originallemma}{Original lemma}{Original lemmas}
\makeatletter
\newlabel{more-algebra-lemma-tensor-NL}{{0FUY}{source}{}{sourcepointer-5}{}}
\newlabel{more-algebra-lemma-tensor-NL@cref}{{[originallemma][0][]0FUY}{[1][0][]source}{}{}{}}
\makeatother
\title{The naive cotangent complex computed from a factorization through a formally smooth morphism}
\author{The Stacks Project Authors}
\date{stacks-project:more-morphisms.tex:local-snapshot-20261011; GNU FDL 1.2 or later}
\begin{document}
\maketitle
\noindent\textit{Editorial note.} Statement span is lines 3979-3989 of sources/stacks-more-morphisms/pinned-original-source.tex; proof span is lines 3991-4015, the complete author proof (23 body lines).\par\medskip
\noindent\textit{Editorial note.} Navigation stubs: 5 cited-but-not-included labels are supplied as non-load-bearing pointers with their pinned tags where the snapshot tag index has them.\par\medskip
\noindent\textit{Editorial note.} extra\_packages is empty on purpose: the retained preamble already loads xy with options and a second declaration clashes in the assembler wrapper.\par\medskip
\noindent\textit{Editorial note.} Text is retained verbatim, including the author's own typos and placeholders.\par\medskip
\noindent\textit{Allowed standard prerequisites (editorial).}\begin{itemize}
\item Explicit prerequisite boundary. The naive cotangent complex, conormal sheaves of immersions, and formally smooth morphisms; the comparison is retained.
\end{itemize}\medskip
\ifdefined\hypertarget\hypertarget{sourcepointer-1}{}\fi
\noindent\textit{Source reference.} Original lemma lemma-NL-formally-smooth: sources/stacks-more-morphisms/pinned-original-source.tex, source lines 3824--3830. Cited by the retained statement or proof but not part of the collected spans; supplied only as a navigation pointer. The argument does not depend on its content.\par\medskip
\ifdefined\hypertarget\hypertarget{sourcepointer-2}{}\fi
\noindent\textit{Source reference.} Original lemma lemma-NL-immersion: sources/stacks-more-morphisms/pinned-original-source.tex, source lines 3903--3907. Cited by the retained statement or proof but not part of the collected spans; supplied only as a navigation pointer. The argument does not depend on its content.\par\medskip
\ifdefined\hypertarget\hypertarget{sourcepointer-3}{}\fi
\noindent\textit{Source reference.} Original lemma lemma-smooth-formally-smooth: sources/stacks-more-morphisms/pinned-original-source.tex, source lines 3190--3197. Cited by the retained statement or proof but not part of the collected spans; supplied only as a navigation pointer. The argument does not depend on its content.\par\medskip
\ifdefined\hypertarget\hypertarget{sourcepointer-4}{}\fi
\noindent\textit{Source reference.} Original lemma 0E1Z: sources/stacks-more-morphisms/pinned-original-source.tex, source lines 4015--4015. Cited by the retained statement or proof but not part of the collected spans; supplied only as a navigation pointer. The argument does not depend on its content.\par\medskip
\ifdefined\hypertarget\hypertarget{sourcepointer-5}{}\fi
\noindent\textit{Source reference.} Original lemma 0FUY: sources/stacks-more-morphisms/pinned-original-source.tex, source lines 4015--4015. Cited by the retained statement or proof but not part of the collected spans; supplied only as a navigation pointer. The argument does not depend on its content.\par\medskip
\begin{lemma}
\label{lemma-get-NL}
Let $f : X \to Y$ be a morphism of schemes which factors
as $f = g \circ i$ with $i$ an immersion and $g : P \to Y$
formally smooth (for example smooth). Then there is a canonical isomorphism
$$
\NL_{X/Y} \cong \left(\mathcal{C}_{X/P} \to i^*\Omega_{P/Y}\right)
$$
in $D(\mathcal{O}_X)$ where the conormal sheaf $\mathcal{C}_{X/P}$
is placed in degree $-1$.
\end{lemma}

\begin{proof}
(For the parenthetical statement see Lemma \ref{lemma-smooth-formally-smooth}.)
By Lemmas \ref{lemma-NL-immersion} and \ref{lemma-NL-formally-smooth} we have
$\NL_{X/P} = \mathcal{C}_{X/P}[1]$ and $\NL_{P/Y} = \Omega_{P/Y}$ with
$\Omega_{P/Y}$ locally projective. This implies that
$i^*\NL_{P/Y} \to i^*\Omega_{P/Y}$ is a quasi-isomorphism too
(small detail omitted; the reason is that $i^*\NL_{P/Y}$ is the
same thing as $\tau_{\geq -1}Li^*\NL_{P/Y}$, see More on Algebra, Lemma
\ref{more-algebra-lemma-tensor-NL}).
Thus the canonical map
$$
i^*\NL_{P/Y} \to \text{Cone}(\NL_{X/Y} \to \NL_{X/P})[-1]
$$
of Modules, Lemma \ref{modules-lemma-exact-sequence-NL-ringed-topoi}
is an isomorphism in $D(\mathcal{O}_X)$ because the cohomology
group $H^{-1}(i^*\NL_{P/Y})$ is zero by what we said above.
In other words, we have a distinguished triangle
$$
i^*\NL_{P/Y} \to \NL_{X/Y} \to \NL_{X/P} \to i^*\NL_{P/Y}[1]
$$
Clearly, this means that $\NL_{X/Y}$ is the cone on the map
$\NL_{X/P}[-1] \to i^*\NL_{P/Y}$ which is equivalent to the
statement of the lemma by our computation of the cohomology
sheaves of these objects in the derived category given above.
\end{proof}

\end{document}
